\section{Discussion}
\label{sec:discussion}
\paragraph{Different measurements support different claims.}
The studies form an evidence ladder rather than a single ranking. Correctness transitions describe practical changes. Explanation and random-source controls narrow their attribution. The factorial contrast provides evidence that answer benefits depend on designed correspondence. Explicit artifacts then expose content deficiencies and conditional continuation outcomes. No one layer establishes all the others, and evidence of supplied-source use does not automatically validate a self-generated staged path.

\paragraph{What the framework adds.}
The framework turns a general recommendation to ``inspect reasoning'' into specific questions: does the source contain the necessary relation; does the mapping retain it; does the schema preserve applicability conditions; and does adaptation satisfy target constraints? These questions can guide source curation, prompt inspection, and error analysis. The empty-schema successes give an objective example: accuracy can credit an answer even when the designed path delivers no schema content. The remaining stage assignments require semantic inspection, while the recorded contracts and this absence check are mechanically recoverable. Their practical value is the distinction between available information and successful use of it.

\paragraph{What remains unmeasured.}
Stage-label agreement, blind reference validation, matched rephrasing controls, unchanged-artifact reruns, and information-access ablations remain outstanding. Consequently, we present a diagnostic protocol with descriptive case-study evidence, not a validated automatic causal attribution system. The all-R GPQA pattern particularly merits separating factual source errors from relational inadequacy and inspection bias. Restricting middle-stage inputs may itself contribute to failure. These questions can refine the framework; the present study does not resolve them by treating the stage outputs as internal cognition.
